%-----------------------------------------------------------------------
% declaration_ai.tex —— 四川大学本科毕业论文 AI 工具使用声明

\chapter*{四川大学本科毕业论文 AI 工具使用声明}
\ctexset{section/format = {\raggedright\zihao{3}\heiti}}

本人声明，在撰写题目为《\varTitle》的本科毕业论文过程中，对 AI 工具的使用情况如下：

\section*{一、使用的 AI 工具及使用目的}

\begin{table}[H]
\centering
\renewcommand{\arraystretch}{1.3}
\begin{tabular}{|c|c|p{6cm}|}
\hline
序号 & AI 工具 & 使用目的 \\
\hline
1  & （填写工具名称） & （如：语言润色、语法纠正、代码生成等） \\
\hline
\multicolumn{3}{c}{} \\
\end{tabular}
\end{table}

\section*{二、使用程度说明}

\begin{table}[H]
\centering
\renewcommand{\arraystretch}{1.3}
\begin{tabular}{|c|c|c|c|c|c|}
\hline
序号 & 章节 & AI 工具 & 使用方式 & 内容占比(\%) & 处理方式 \\
\hline
1 & （章节名） & （工具名） & （简要描述） & （百分比） & （如：重新组织语言、检查语法错误等） \\
\hline
\multicolumn{6}{c}{} \\
\end{tabular}
\end{table}

在论文的（具体章节）使用了上述 AI 工具。使用方式为（简要说明）。在这些部分，AI 工具提供的内容占比约为 xx\textbf{\%}。

本人对 AI 工具提供的内容进行了（处理方式），以确保论文内容符合学术要求及本人的研究思路。

\section*{三、原创性与责任声明}

本人确认，尽管使用了上述 AI 工具，但论文整体的研究思路、核心观点、论证过程及最终结论均为本人独立思考与研究的成果。论文中引用 AI 工具生成内容之处，均已按照学校规定的学术规范进行了明确标注。

若论文因 AI 工具使用不当或未按要求标注等原因出现学术不端问题，本人愿意承担全部责任。

\vspace{20pt}

作者签名：\rule{3cm}{1.2cm}
\hfill
导师签名：\rule{3cm}{1.2cm}

\vspace{0.5cm}

\begin{flushright}
    \makebox[3.5cm][r]{xxxx年xx月xx日}
\end{flushright}
